% !TEX root = ../notes.tex

\documentclass[../notes.tex]{subfiles}

\begin{document}

\section{April 1}
Here we go.

\subsection{Real Forms of Classical Types}
Let's list out some real forms we know.
\begin{exe}
	We list some real forms of $\mf{sl}_n(\CC)$.
\end{exe}
\begin{proof}
	Here are our forms.
	\begin{itemize}
		\item We have the split real form $\mf{sl}_n(\RR)$
		\item One can directly compute that the compact real form is $\mf{su}_n(\RR)$: one can calculate that $\tau$ is given by the Lie algebra involution $\sigma_\tau(X)=-X^\dagger$.
		\item Because $\op{Aut}A_{n-1}$ is $\ZZ/2\ZZ$ (for $n\ge3$, given by the ``flip'' of the Dynkin diagram), there is a unique non-inner quasisplit real form. It should be given by the $s_0$ which sends $e_j\mapsto e_{n-j}$ and $f_j\mapsto f_{n-j}$ and $h_j\mapsto h_{n-j}$. One can check that this automorphism is actually given by $s_0(A)\coloneqq-JA^\intercal J^{-1}$, where $J$ is the antidiagonal matrix
		\[J=\begin{bmatrix}
			&&&& -1 \\
			&&& 1 \\
			&& -1 \\
			& 1 \\
			\iddots
		\end{bmatrix}.\]
		Thus, this quasisplit form is $\mf{su}(p,p)$ if $n=2p$, and it is $\mf{su}(p,p+1)$ if $n=2p+1$. Note that comparing $\tau$ and $s_0$ shows that the compact form and this quasisplit form are in the same inner class!
		\item In general, one has the real forms $\mf{su}(p,q)$ where $n=p+q$.
		\qedhere
	\end{itemize}
\end{proof}
\begin{remark}
	For $n=2$, one has that $\mf{su}(1,1)=\mf{sl}(2)$ because the Dynkin diagram has no nontrivial automorphisms.
\end{remark}
\begin{exe}
	We list some real forms of $\mf{so}(2n+1,\CC)$.
\end{exe}
\begin{proof}
	Here are our forms. Note that all forms are inner because there are no nontrivial automorphism of the Dynkin diagram.
	\begin{itemize}
		\item One has the compact form $\mf{so}(2n+1)$.
		\item There is also the split form $\mf{so}(n,n+1)$, which one sees from the Serre presentation (or alternatively, the standard form of $\mf{so}(2n+1)$ over $\CC$).
		\item In general, there are forms $\mf{so}(p,q)$ when $p+q=2n+1$.
	\end{itemize}
	Note that when $n=1$, we have $\mf{so}(1,2)=\mf{sl}(2,\RR)$.
\end{proof}
\begin{exe}
	We list some real forms of $\mf{sp}(2n,\CC)$.
\end{exe}
\begin{proof}
	Note that all forms are inner because the Dynkin diagram has no nontrivial automorphisms. There is the compact form $\mf{su}(n,\HH)$ and the split form $\mf{sp}(2n,\RR)$. Note that when $n=2$, this should agree with $\mf{so}(2n+1)$, so we note that $\mf{su}(2,\HH)=\mf{so}(5)$ and $\mf{sp}(4,\RR)=\mf{so}(3,2)$ due to the uniqueness of the split and compact forms.
\end{proof}
\begin{exe}
	We list some real forms of $\mf{so}(2n,\CC)$.
\end{exe}
\begin{proof}
	There is a compact form. When $n\ge3$, there are two automorphisms of the Dynkin diagram; even when $n=2$, there are only two conjugacy classes of automorphisms. Thus, there is a compact form $\mf{so}(2n)$, a split form $\mf{so}(n,n)$ (as in type $B$), and there is a unique non-inner quasisplit form. By noting that the nontrivial conjugacy class may as well be given by $e_n\mapsto-e_n$ and $e_i\mapsto e_i$ for $i<n$, we may take the quadratic form of the quasisplit form to come from the symmetric matrix
	\[\begin{bmatrix}
		0 & 0 & 1_{n-1} & 0 \\
		0 & 1 & 0 & 0 \\
		1_{n-1} & 0 & 0 & 0 \\
		0 & 0 & 0 & 1
	\end{bmatrix},\]
	which has signature $(n+1,n-1)$.
\end{proof}
\begin{example}
	For $D_2$, one can see that the quasisplit inner form recovers $\mf{sl}(2,\CC)$ because $D_2=A_2$. Thus, $\mf{sl}(2,\CC)\cong\mf{so}(3,1)$ (as real Lie algebras). Note that the corresponding real points of the group $\op{SO}(3,1)$ has two connected components! Indeed, $\op{SO}(3,1)$ acts naturally on the cone cut out by $t^2-x_1^2-x_2^2-x_3^2=0$, and there is a real Lie subgroup preserving the points with $t>0$.
\end{example}
\begin{example}
	For $n=3$, one has $D_3=A_3$, and one finds that $\mf{so}(6)\cong\mf{su}(4)$ and $\mf{so}(3,3)\cong\mf{sl}(4,\RR)$ and $\mf{so}(4,2)\cong\mf{su}(2,2)$, again by matching up compact, slit, and quasisplit forms.
\end{example}
One can say something similar for exceptional types, but we will return to them later.

\subsection{Counting from the Compact Form}
For real groups, it turns out to be more convenient to use Galois cohomology to count ``from'' the compact form instead of the split form.
\begin{proposition} \label{prop:form-by-compact}
	Fix a complex semisimple Lie algebra $\mf g$. Then the real forms of $\mf g$ are in bijection with conjugacy classes of involutions in $\op{Aut}\mf g^c$.
\end{proposition}
\begin{proof}
	Here, we start with the antilinear Cartan involution $\omega$ coming from the compact Cartan involution. We now run the same argument as in \Cref{prop:classify-ss-forms}. Now, another real structure on $\mf g$ can be described by an antilinear involution, which we may take to have the form $\sigma\coloneqq\omega\circ g$, where $g\in\op{Aut}(\mf g)$ is chosen with $\omega(g)g=1$.

	To continue, we claim that
	\[\omega(x)\stackrel?=x^{-\dagger},\]
	where $(\cdot)^\dagger$ is induced from the Killing form on $\mf g^c$, extended to a Hermitian form. Indeed, to show the claim, it is enough to show that $\omega(X)=-X^\dagger$ on the Lie algebra; by splitting up into Hermitian and skew-Hermitian parts, it is enough to note that $\omega(X)=X$ if $X$ is unitary for the Hermitian form on $\mf g$ coming from the Killing form on $\mf g^c$ (because $\omega$ is negative-definite for the Killing form on $\mf g^c$).

	The moral is that $\omega(g)g=1$ implies that
	\[g^\dagger g=1,\]
	so $g$ is self-adjoint, so it is self-adjoint with real eigenvalues. Thus, we may decompose
	\[\mf g=\bigoplus_{\gamma\in\RR}\mf g[\gamma],\]
	where $\mf g[\gamma]$ is the $\gamma$-eigenspace of $\mf g$. Because $g$ is an automorphism of $\mf g$, we see that $[\mf g[\gamma],\mf g[\gamma']]=\mf g[\gamma\gamma']$.

	We now modify $g$ for convenience. For any $t$, we define $\left|g\right|^t$ to act by $\left|\gamma\right|^t$ on $\mf g[\gamma]$; note that $\left|g\right|^t$ continues to be an automorphism of $\mf g$, which one can see by an eigenvalue decomposition. Now, define $\theta\coloneqq g\left|g\right|^{-1}$, so it follows that $\theta^2=1$ by looking at eigenspace actions. Because $g\omega=\omega g^{-1}$ from earlier, we see that $\theta\omega=\omega\theta^{-1}=\omega\theta$. Additionally, $\theta=\left|g\right|^{-1/2}g\left|g\right|^{1/2}$, so $\theta$ and $g$ are conjugate, so they define the same real form.

	The previous paragraph allows us to replace $g$ with $\theta$, which now squares to $1$ and commutes with $\omega$; commuting with $\omega$ is enough to view $\theta$ as an automorphism of $\mf g^c$, so we have at least described how to take a real form and produce some $\theta$. For such $\theta$, we see that another $\theta'$ gives the same real form when $\theta'=x\theta\omega(x)^{-1}$ for some automorphism $x$. To this end, we write $z\coloneqq\omega(x)^{-1}x$, which is $x^\dagger x$ and thus positive-definite. By construction, we see that $\omega(z)=z^{-1}$ because $\omega$ is antilinear with $\omega^2=1$, and the construction gives $\theta z=z^{-1}\theta$.

	Accordingly, we normalize $x$ with $y\coloneqq xz^{-1/2}$ so that $\omega(y)=\omega(x)z^{1/2}=xz^{-1/2}=y$. Similarly, we may calculate
	\begin{align*}
		\theta' &= x\theta\omega(x)^{-1} \\
		&= x\theta zx^{-1} \\
		&= xz^{-1/2}\theta z^{1/2}x^{-1} \\
		&= y\theta y^{-1},
	\end{align*}
	so $\theta$ is actually conjugate by this nice element $y$, which is $\omega$-invariant!
\end{proof}
\begin{remark}
	Given some $\theta$, the proof shows that our involution for \Cref{prop:classify-ss-forms} is given by $\omega_\theta\coloneqq\omega\circ\theta$.
\end{remark}
\begin{remark} \label{rem:get-g-theta}
	Now, given $\theta$, one can split $\mf g$ into the complex subspaces $\mf k\oplus\mf p$, which are the $(+1)$- and $(-1)$-eigenspaces, respectively. We can see that $\mf k$ is a Lie subalgebra, and the definiteness of the Killing form implies that $\mf k$ is reductive by Cartan's criterion. Now, this decomposition commutes with $\omega$, so we may write
	\[\mf g^c=\mf k^c\oplus\mf p^c.\]
	The corresponding real form $\mf g_\theta$ can thus be seen to be the $\omega_\theta$-conjugate-invariants. Thus, this consists of the invariants in $\mf k^c$ (because here $\theta$ acts by $1$) and consists of $i\mf p^c$ (because here $\theta$ acts by $-1$), so $\mf g_\theta=\mf k^c\oplus i\mf p^c$ by gluing these together.
\end{remark}

\subsection{Vogan Diagrams}
We will want to use \Cref{prop:form-by-compact} to classify real forms by Dynkin diagrams. Here is a quick coherence check, which allows us to do a root decomposition of $\mf k$ and $\mf g$ simultaneously.
\begin{proposition} \label{prop:get-cartan-g-theta}
	Fix a real form $\mf g_\theta$ of a complex semisimple Lie algebra $\mf g$, and set $\mf k\coloneqq\mf g^\theta$. Then there is a $\theta$-stable Cartan subalgebra $\mf h$ of $\mf g$ for which $\mf h\cap\mf k$ is a Cartan subalgebra of $\mf k$.
\end{proposition}
\begin{proof}
	Fix everything as in \Cref{rem:get-g-theta}. Choose a generic $t\in\mf k^c$, which is known to be regular semisimple, and we let $\mf h_+^c$ be its centralizer in $\mf k^c$. The complexification is called $\mf h_+$, and because $t$ is regular, we see that $\mf h_+\subseteq\mf h$ is a Cartan subalgebra of $\mf k$.

	We now extend $\mf h_+$ to $\mf g$. Indeed, let $\mf h_-^c$ be a maximal subspace of $\mf p^c$ so that
	\[\mf h^c\coloneqq\mf h_+^c\oplus\mf h_-^c\]
	is an abelian subalgebra of $\mf g^c$. We claim that the complexification $\mf h\subseteq\mf g$ is a Cartan subalgebra, which will complete the proof by construction (note that $\mf h\cap\mf k=\mf h_+$ because $\mf p\cap\mf k=0$). Here are our checks.
	\begin{itemize}
		\item Semisimple: note that $\mf h^c$ consists of semisimple elements because these elements are anti-Hermitian operators on $\mf g^c$.
		\item Maximal: suppose that $z$ commutes with $\mf h$. Then the decomposition $\mf g=\mf k\oplus\mf p$ allows us to decompose $z=z_++z_-$. Further, because $\theta$ commutes with $\omega$, we see that $z_+$ must commute with $\mf h$, so $z_+\in\mf h_+$ automatically. Thus, we may assume that $z_+=0$. Then we write $z_-=x+iy$ where $x,y\in\mf p^c$, and both $x$ and $y$ continue to commute with $\mf h$ by taking real and complex parts, so $x,y\in\mf h_-$ by maximality. Thus, $z_-\in\mf h_-$ as well.
		\qedhere
	\end{itemize}
\end{proof}
\begin{remark}
	With the above construction, we further claim that $\mf h_-$ does not have any coroots of $\mf g$. Indeed, suppose that we have a coroot $\alpha^\lor\in\mf h_-$. This means that $\theta(\alpha^\lor)=-\alpha^\lor$ because $\mf h_-\subseteq\mf p$, so $\theta(e_\alpha)=e_{-\alpha}$ and $\theta(e_{-\alpha})=e_\alpha$ follows. Now, set $x\coloneqq e_\alpha+e_{-\alpha}$, which we see is $\theta$-stable, so $x\in\mf k$. Now, on one hand, $[\mf h_+,x]=0$ because $\alpha$ vanishes on $\mf h_+$, but $\mf h_+\subseteq\mf k$ is a maximal abelian subalgebra, so it must contain $x$! On the other hand, the root decomposition means that $x$ is orthogonal to $\mf h_+$, so $x=0$, which is a contradiction.
\end{remark}
We now do our root decompositions. Fix a Cartan as in \Cref{prop:get-cartan-g-theta}. Then we may choose a generic $t\in\mf h_+$, and we see that $t$ is regular in $\mf g$ because $\mf h_+$ has all the roots. Additionally, generic $t$ will have $\op{Re}(t,\alpha^\lor)\ne0$ for all coroots $\alpha^\lor$, which we use to produce a polarization. By construction, we finally see that $\theta$ preserves the set of positive roots (because $t$ is $\theta$-invariant!), so $\theta$ preserves the set of simple roots, so $\theta$ induces an automorphism of the Dynkin diagram. Let $\{\alpha_i\}$ be the collection of simple roots. There are two cases.
\begin{itemize}
	\item If $\theta(i)=i$, then $\theta(e_i)=\pm e_i$ and $\theta(f_i)=\pm f_i$ and $\theta(h_i)=\pm h_i$. Indeed, this is required because $\theta$ preserves this root space, and $\theta^2=1$.
	\item Otherwise, we can have $\theta(i)\ne i$, but then $\theta(h_i)=h_{\theta(i)}$, and we can normalize everything so that $\theta(e_i)=e_{\theta(i)}$ and $\theta(f_i)=f_{\theta(i)}$.
\end{itemize}
The moral is that this choice of $\theta$ and $t$ can be labeled by a Vogan diagram, which is a decoration of the Dynkin diagram as follows: we choose an automorphism $\theta$ of the Dynkin diagram, and we must color the $\theta$-invariant vertices either white or black (corresponding to the sign $+$ or $-$ in the equation $\theta(e_i)=\pm e_i$, respectively).
\begin{remark}
	We have shown that every real form $\mf g_\theta$ is labeled by at least one Vogan diagram coming from $\theta$ and a choice of polarization. But more than one Vogan diagram may produce the same real form!
\end{remark}
\begin{example}
	For $A_1$, there are only two Vogan diagrams: we have a single white dot, which is $\mf{sl}(2,\RR)$; and we have a single black dot, which is $\mf{su}(2)$.
\end{example}

\subsection{Real Forms of Type \texorpdfstring{$A$}{ A}}
Let's use \Cref{prop:form-by-compact} for fun and profit.
\begin{exe} \label{exe:real-form-a}
	We classify the real forms of type $A_{n-1}$.
\end{exe}
\begin{proof}
	Here, $\mf g=\mf{sl}_n(\CC)$, so $\mf g^c=\mf{su}(n)$, which has automorphism group $\ZZ/2\ZZ\times\op{PU}(n)$; namely, these are the automorphisms of $\mf g$ which preserve $\mf g^c$. There are thus two cases.
	\begin{itemize}
		\item We have a compact inner class, which comes from those $\theta\in\op{PU}(n)$. Up to conjugacy, we may move such matrices to have $p$ total $(+1)$s and $q$ total $(-1)$s, where $p$ and $q$ are nonnegative integers with $n=p+q$. We can now directly calculate that $\mf g_\theta=\mf{su}(p,q)$. It turns out that $\mf k=\mf{sl}(p)\oplus\mf{sl}(q)\oplus\RR$.

		\item We have a split inner class, which corresponds to $\theta$ having nontrivial image in $\ZZ/2\ZZ$, so $\theta$ acts by the flip on the Dynkin diagram. If $n$ is odd, then there is only one Vogan diagram (because $\theta$ fixes nothing), so there is only the split form.

		Otherwise, if $n$ is even, then there are two Vogan diagrams (because $\theta$ fixes the middle vertex), so we have at most two real forms. Here, one calculates that our real forms are the split form $\mf{sl}(n,\RR)$ with $\mf k=\mf{so}(n)$ (for black vertex) and the additional form $\{X\in\mf{su}(n/2,\HH):\Re\tr X=0\}$ with $\mf k=\mf{sp}(n)$ (with white vertex).
		\qedhere
	\end{itemize}
\end{proof}
\begin{remark}
	Already we can see that there are far more Vogan diagrams than real forms in the case where $\theta$ acts trivially.
\end{remark}
\begin{example} \label{ex:compute-form-from-vogan}
	Consider the following Vogan diagram
	% https://q.uiver.app/#q=WzAsNixbMCwwLCJcXGNpcmMiXSxbMSwwLCJcXGJ1bGxldCJdLFsyLDAsIlxcY2lyYyJdLFszLDAsIlxcY2lyYyJdLFs0LDAsIlxcYnVsbGV0Il0sWzUsMCwiXFxjaXJjIl0sWzAsMSwiIiwwLHsic3R5bGUiOnsiaGVhZCI6eyJuYW1lIjoibm9uZSJ9fX1dLFsxLDIsIiIsMCx7InN0eWxlIjp7ImhlYWQiOnsibmFtZSI6Im5vbmUifX19XSxbMiwzLCIiLDAseyJzdHlsZSI6eyJoZWFkIjp7Im5hbWUiOiJub25lIn19fV0sWzMsNCwiIiwwLHsic3R5bGUiOnsiaGVhZCI6eyJuYW1lIjoibm9uZSJ9fX1dLFs0LDUsIiIsMCx7InN0eWxlIjp7ImhlYWQiOnsibmFtZSI6Im5vbmUifX19XV0=&macro_url=https%3A%2F%2Fraw.githubusercontent.com%2FdFoiler%2Fnotes%2Fmaster%2Fnir.tex
	\[\begin{tikzcd}[cramped]
		\circ & \bullet & \circ & \circ & \bullet & \circ
		\arrow[no head, from=1-1, to=1-2]
		\arrow[no head, from=1-2, to=1-3]
		\arrow[no head, from=1-3, to=1-4]
		\arrow[no head, from=1-4, to=1-5]
		\arrow[no head, from=1-5, to=1-6]
	\end{tikzcd}\]
	which gives a form of $\mf{sl}_7(\CC)$. Here, we see that $\theta(E_{12})=E_{12}$, $\theta(E_{23})=-E_{23}$, and so on. Thus, we are looking at the Lie algebra of the special unitary group preserving the bilinear form given by $\theta=\op{diag}(1,1,-1,-1,-1,1,1)$, which we see by examining sign changes. Thus, we have the real form $\mf{su}(4,3)$! A similar argument shows that each black vertex corresponds to a sign change (whereas a white vertex means staying the same sign), allowing us to calculate the real form for any such Vogan diagram.
\end{example}

\end{document}